\chapter{PDDL Use}
\label{ch:pddl}

The BDI loop is complemented by a symbolic planning layer reserved for one
sub-problem: clearing a pushable crate that blocks a corridor. PDDL(Planning Domain Definition Language) 
is an exact sybolic planner, useless for reward economics but ideal for the discrete, combinatorial
question ``what sequence of moves and pushes relocates this crate out of the way?'',
where an ad-hoc heuristic would be brittle.

\section{Role and scope}
\label{sec:role_scope}
The design choice is to invoke it as narrowly as possible. Ordinary tile-by-tile
navigation stays with A\textsuperscript{*}, which is fast and runs constantly; a
general-purpose solver called on every query would be far too slow for a
\SI{60}{\hertz} control loop. Restricting PDDL to occasional obstacle resolution keeps
it off the main loop while still guaranteeing a correct plan where local heuristics fail.
\\
The planner is reached through the \texttt{pddl-client} interface. For testing and
development we point it at a \emph{local} service on \texttt{http://localhost:5001}
(set via the \texttt{PAAS\_HOST} variable) rather than the public hosted endpoint, for
performance: the online solver takes \SIrange{1}{3}{\second} per query plus network
latency, whereas a local instance answers far faster. Since crate clearing suspends the
agent's active plan, cutting that latency directly shortens the time it stands idle
before an obstacle.

\section{The domain}
\label{sec:domain}
The domain, \texttt{deliveroo-crates}, is a compact STRIPS(STanford Research Institute Problem Solver) 
model with two families of actions:

\begin{itemize}
  \item \textbf{Movement} (\texttt{move-right/left/up/down}): the agent moves to
        an adjacent tile, provided that it's marked \texttt{clear}, 
        meaning it is currently known to be free of crates and other agents.
  \item \textbf{Pushing} (\texttt{push-right/left/up/down}): when the agent, a
        crate, and a destination tile are aligned in the same direction, the agent
        advances onto the crate's tile and push the crate one step forward in a clear and allowed direction(crate tile, crate-move-capable flag).
        The effects flip the clear predicates accordingly.
\end{itemize}

Directionality is encoded as explicit adjacency predicates
(right/left/up/down) rather than coordinates, so the planner reasons over
a symbolic graph.  Because these predicates are emitted only when the move-legality check passes, the
one-way gate rules are respected automatically.

\section{The problem}
\label{sec:problem}
Each time a crate must be cleared, a fresh problem is compiled from the agent's
current beliefs. The compiler scans the walkable tiles and produces four kinds of
facts:

\begin{itemize}
  \item \textbf{Tiles and connectivity:} one object per walkable tile, together with
        the directional adjacency facts, obtained from the same move-legality check
        the agent uses for ordinary navigation.
  \item \textbf{Crate-capable tiles:} a \texttt{crate-move-capable} flag on the tiles
        that can physically hold a crate (the crate-spawn and crate-move types).
  \item \textbf{Free tiles:} a \texttt{clear} flag on every tile not occupied by a
        crate or a peer agent (both treated as obstacles).
  \item \textbf{Objects:} the agent's position, the target crate, and every other
        crate, each as a distinct object.
\end{itemize}

A design choice that has been applied is that only \emph{physical} obstacles mark a tile
as non-clear. \emph{Transient} goal blocks, such as a delivery zone that is
momentarily unreachable, are left out, because including them would label otherwise
reachable tiles as blocked and could lead the planner to wrongly conclude that no
solution exists.

The goal is a single fact: the blocking crate relocated onto a chosen clearance
tile, that is, an empty crate-capable tile out of the agent's way. The plan the
solver returns, a sequence of moves and pushes, is then translated back into grid
steps and handed to the executor as an interruptible plan.

\section{Focus on crates and the tiered strategy}

Crates are the only dynamic, \emph{relocatable} obstacle in the Deliveroo game, which is why
they receive dedicated planning. The agent first picks a sensible destination for
the blocking crate: a nearby clear, crate-capable tile, reachable via a push whose
``push-from'' side the agent can actually stand on (found by a short breadth-first
search). Resolution then proceeds in tiers, cheapest first:

\begin{center}
\begin{verbatim}
  A* path blocked by crate
        |
        v
  pick destination tile  --(none)-->  block goal temporarily
        |
        v
  local push solver  --(success)-->  execute push          (0 ms)
        |
     (fails)
        v
  PDDL solver @ :5001  --(success)-->  execute plan
        |
     (fails)
        v
  record cooldown + block goal
\end{verbatim}
\end{center}

A purely local, single-push computation is tried first because it is instantaneous
and covers the common case where the crate is bypassable using and adjacent tile.
When that fails, typically because the push
requires manoeuvring around further crates, the agent falls back to the
PDDL solver, which can chain multiple pushes and moves. Failed solves are recorded
with a 30\,s cooldown so the agent does not repeatedly hammer the planner on an
unsolvable blockage; in that case the blocked target is shelved temporarily and the
deliberation layer simply chooses something else to do and so picks a different goal.
